% Abstract in English.
\setlength{\absparsep}{18pt}
\begin{resumo}[Abstract]
\begin{otherlanguage*}{english}


\textbf{Context:} Data-intensive systems are fundamental for processing, storing, moving, and analyzing large volumes of data in modern software infrastructures. These systems involve architectural decisions related to processing paradigms, data ingestion, scalability, resource usage, and fault tolerance, which directly impact attributes such as \textit{throughput}, latency, stability, and computational cost. \textbf{Problem:} Despite the existence of different techniques, \textit{frameworks}, and optimization strategies for data-intensive systems, the empirical evidence available in the literature remains fragmented, heterogeneous, and often difficult to compare. This limitation makes evidence-based architectural decision-making difficult, especially in scenarios that require choosing between batch processing, stream processing, and different execution configurations. \textbf{Solution:} This research proposes to investigate a quantitative approach to support architectural decision-making in data-intensive systems by integrating evidence from the literature, experimental results, and observable data characteristics, such as volume, entropy, redundancy, cardinality, and variability. \textbf{Method:} The research was organized into three complementary stages: a Systematic Literature Review, a controlled experimental campaign, and the proposal of a quantitative model to support architectural decision-making. The systematic review identified techniques, metrics, gaps, and evidence related to the design of data-intensive systems. The experimental campaign compared \textit{batch} and \textit{streaming} processing in an environment based on Apache Spark, Spark Structured Streaming, Apache Kafka, and Docker, using the NYC Taxi dataset. Different data volumes, event rates, numbers of \textit{workers}, and \textit{trigger} intervals were evaluated, considering metrics such as \textit{throughput}, latency, CPU usage, memory consumption, \textit{backpressure}, and scalability. \textbf{Summary of results:} The systematic review identified 979 records, of which ten primary studies met the inclusion and quality criteria. The exploratory quantitative synthesis indicated positive effects of architectural optimizations, with an average response ratio of approximately $2.7\times$ compared to baseline configurations. The preliminary experimental results indicated that \textit{batch} processing achieved higher \textit{throughput} in the evaluated scenarios, while \textit{streaming} processing showed greater sensitivity to event rate, \textit{trigger} interval, and micro-batch coordination costs. \textbf{Contribution:} This research is expected to contribute to architectural decision-making in data-intensive systems through an approach based on empirical evidence, experimental results, and informational characteristics of data, supporting software architects and data engineers in the preliminary comparison of architectural alternatives before the complete implementation of a solution.

\textbf{Keywords:} Data-Intensive Systems. Software Architecture. Architectural Decision-Making. Stream Processing. Information Theory.



\end{otherlanguage*}
\end{resumo}
